% =============================================================
% ПОРОЖДЁННЫЙ ФАЙЛ — НЕ РЕДАКТИРОВАТЬ ВРУЧНУЮ.
%
% Источник:  Build/css/enf-tokens.css
% Генератор: Scripts/gen-latex-colors.mjs
% Пересоздать: node Scripts/gen-latex-colors.mjs
%
% Берётся светлая тема: печать идёт по белой бумаге.
% =============================================================

\usepackage{xcolor}

\definecolor{enfvariable}{HTML}{1D5B85}  % объект, над которым действуем
\definecolor{enffunction}{HTML}{8A5214}  % действие, преобразование
\definecolor{enfparameter}{HTML}{3E7D3A}  % то, что настраивается
\definecolor{enfoperator}{HTML}{A32EAD}  % агрегат, свёртка
\definecolor{enftarget}{HTML}{7A0A2C}  % цель: награда, потери, оптимум
\definecolor{enfneutral}{HTML}{414A54}  % служебное обозначение

% Макросы ролей. Имена совпадают с макросами MathJax из хранилища,
% поэтому один и тот же исходник собирается и там, и здесь.

\newcommand{\enfVar}[1]{\textcolor{enfvariable}{#1}}
\newcommand{\enfFun}[1]{\textcolor{enffunction}{#1}}
\newcommand{\enfPar}[1]{\textcolor{enfparameter}{#1}}
\newcommand{\enfOp}[1]{\textcolor{enfoperator}{#1}}
\newcommand{\enfTgt}[1]{\textcolor{enftarget}{#1}}
\newcommand{\enfNeu}[1]{\textcolor{enfneutral}{#1}}

